\documentclass{article}
\usepackage{listings}
\usepackage{xcolor}
\usepackage{graphicx}
\usepackage{float} % For [H] to work

\title{OS\ Project}
\author{22P9326 Mashaal Ishtiaq}


% Define custom colors for code
\definecolor{codebackground}{rgb}{0.95, 0.95, 0.95}
\definecolor{keywords}{rgb}{0.0, 0.0, 1.0}
\definecolor{comments}{rgb}{0.0, 0.5, 0.0}

\lstset{
  backgroundcolor=\color{codebackground},
  basicstyle=\ttfamily\footnotesize,
  keywordstyle=\color{keywords},
  commentstyle=\color{comments},
  frame=single,
  breaklines=true,
  columns=flexible,
  keepspaces=true
}

\begin{document}

\maketitle

\section*{Task 1: CPU Scheduling Implementation}

\begin{lstlisting}[language=C++]
#include <iostream>
#include <fstream>
#include <vector>
#include <string>
#include <queue>
#include <algorithm>
#include <sstream> // Include this for istringstream

using namespace std;

// Structure to store task information
struct Task
{
    string name;
    int priority;
    int burstTime;
};

// Function to read tasks from the file
vector<Task> readTasks(const string &filename)
{
    vector<Task> tasks;
    ifstream file(filename);
    if (!file.is_open())
    {
        cerr << "Failed to open file: " << filename << endl;
        return tasks;
    }

    string line;
    while (getline(file, line))
    {
        size_t pos1 = line.find(',');
        size_t pos2 = line.find(',', pos1 + 1);

        string name;
        int priority, burstTime;

        if (pos1 != string::npos && pos2 != string::npos)
        {
            // For comma-separated values
            name = line.substr(0, pos1);
            priority = stoi(line.substr(pos1 + 1, pos2 - pos1 - 1));
            burstTime = stoi(line.substr(pos2 + 1));
        }
        else
        {
            // For space-separated values
            istringstream iss(line); // Use istringstream to parse space-separated values
            iss >> name >> priority >> burstTime;
        }

        tasks.push_back({name, priority, burstTime});
    }

    file.close();
    return tasks;
}

//--------------------FCFS--------------------
void fcfs(const vector<Task> &tasks)
{
    cout << "\n-------FCFS Scheduling------\n";
    int currentTime = 0;
    cout << "\nTask Name\tBurst Time\tWaiting Time\n";

    for (const auto &task : tasks)
    {
        cout << task.name << "\t\t" << task.burstTime << "\t\t" << currentTime << endl;
        currentTime += task.burstTime;
    }
}

//--------------------SJF--------------------
void sjf(vector<Task> tasks)
{
    cout << "\n-------SJF Scheduling------\n";
    sort(tasks.begin(), tasks.end(), [](const Task &a, const Task &b)
         { return a.burstTime < b.burstTime; });

    int currentTime = 0;
    cout << "\nTask Name\tBurst Time\tWaiting Time\n";

    for (const auto &task : tasks)
    {
        cout << task.name << "\t\t" << task.burstTime << "\t\t" << currentTime << endl;
        currentTime += task.burstTime;
    }
}

//--------------------Priority--------------------
void priorityScheduling(vector<Task> tasks)
{
    cout << "\n-------Priority Scheduling------\n";
    sort(tasks.begin(), tasks.end(), [](const Task &a, const Task &b)
         { return a.priority > b.priority; });

    int currentTime = 0;
    cout << "\nTask Name\tBurst Time\tWaiting Time\n";

    for (const auto &task : tasks)
    {
        cout << task.name << "\t\t" << task.burstTime << "\t\t" << currentTime << endl;
        currentTime += task.burstTime;
    }
}

//--------------------RR--------------------
void roundRobin(const vector<Task> &tasks, int quantum)
{
    cout << "\n-------Round Robin Scheduling (Quantum = " << quantum << ")------\n";
    queue<Task> taskQueue;
    for (const auto &task : tasks)
    {
        taskQueue.push(task);
    }

    int currentTime = 0;
    cout << "\nTask Name\tBurst Time\tCompletion Time\n";

    while (!taskQueue.empty())
    {
        Task currentTask = taskQueue.front();
        taskQueue.pop();

        cout << currentTask.name << "\t\t" << currentTask.burstTime << "\t\t" << currentTime;
        if (currentTask.burstTime <= quantum)
        {
            currentTime += currentTask.burstTime;
            cout << endl;
        }
        else
        {
            currentTime += quantum;
            currentTask.burstTime -= quantum;
            cout << " (Remaining: " << currentTask.burstTime << ")\n";
            taskQueue.push(currentTask);
        }
    }
}

//------------------- Priority with Round Robin Scheduling---------------------------
void priorityWithRoundRobin(vector<Task> tasks, int quantum)
{
    cout << "\n-------Priority with Round Robin Scheduling (Quantum = " << quantum << ")------\n";
    sort(tasks.begin(), tasks.end(), [](const Task &a, const Task &b)
         { return a.priority > b.priority; });

    queue<Task> taskQueue;
    for (const auto &task : tasks)
    {
        taskQueue.push(task);
    }

    int currentTime = 0;
    cout << "\nTask Name\tBurst Time\tCompletion Time\n";

    while (!taskQueue.empty())
    {
        Task currentTask = taskQueue.front();
        taskQueue.pop();

        cout << currentTask.name << "\t\t" << currentTask.burstTime << "\t\t" << currentTime;
        if (currentTask.burstTime <= quantum)
        {
            currentTime += currentTask.burstTime;
            cout << endl;
        }
        else
        {
            currentTime += quantum;
            currentTask.burstTime -= quantum;
            cout << " (Remaining: " << currentTask.burstTime << ")\n";
            taskQueue.push(currentTask);
        }
    }
}

//--------------------Main--------------------
int main()
{
    const string filename = "schedule.txt";
    vector<Task> tasks = readTasks(filename);

    if (tasks.empty())
    {
        cerr << "No tasks were read. Please check the file format or path." << endl;
        return 1;
    }

    cout << "Tasks successfully read from file:\n";
    cout << "\n-------Initial Task Details------\n";
    cout << "\nTask Name\tPriority\tBurst Time\n";
    for (const auto &task : tasks)
    {
        cout << task.name << "\t\t" << task.priority << "\t\t" << task.burstTime << endl;
    }

    fcfs(tasks);
    sjf(tasks);
    priorityScheduling(tasks);
    roundRobin(tasks, 4);             // Quantum = 4
    priorityWithRoundRobin(tasks, 4); // Quantum = 4

    return 0;
}

\end{lstlisting}
\begin{figure}
    \centering
    \includegraphics[width=1\linewidth]{1.png}
    \label{fig:enter-label}
\end{figure}

\begin{figure}
    \centering
    \includegraphics[width=1\linewidth]{2.png}
    \label{fig:enter-label}
\end{figure}

\begin{figure}
    \centering
    \includegraphics[width=1\linewidth]{3.png}
    \label{fig:enter-label}
\end{figure}

\begin{figure}
    \centering
    \includegraphics[width=1\linewidth]{4.png}
    \label{fig:enter-label}
\end{figure}

\begin{figure}
    \centering
    \includegraphics[width=1\linewidth]{5.png}

    \label{fig:enter-label}
\end{figure}
\newpage
\section*{Task 2 }
\subsection*{Part 1: Local System Socket Programming}
\subsection*{Server Code}
\begin{lstlisting}[language=C++]
#include <iostream>
#include <vector>
#include <thread>
#include <cstring>
#include <algorithm>
#include <sys/socket.h>
#include <arpa/inet.h>
#include <unistd.h>
using namespace std;

#define PORT 8094
#define BUFFER_SIZE 1024

vector<int> client_sockets;  // List of connected client sockets

// Broadcast a message to all connected clients
void broadcast_message(const string &message) {
    for (int sock : client_sockets) {
        int result = send(sock, message.c_str(), message.size(), 0);
        if (result == -1) {
            cerr << "Failed to send message to client socket: " << sock << "\n";
        } else {
            cout << "Message sent to client socket: " << sock << "\n";
        }
    }
}

// Handle messages from a specific client
void handle_client(int client_socket) {
    char buffer[BUFFER_SIZE];
    while (true) {
        memset(buffer, 0, BUFFER_SIZE);
        int bytes_read = recv(client_socket, buffer, BUFFER_SIZE, 0);
        if (bytes_read <= 0) {
            cerr << "Client disconnected.\n";
            close(client_socket);
            client_sockets.erase(
                remove(client_sockets.begin(), client_sockets.end(), client_socket),
                client_sockets.end()
            );
            break;
        }

        cout << "Received from client: " << buffer << "\n";
        broadcast_message(buffer);
    }
}

int main() {
    int server_socket, client_socket;
    struct sockaddr_in server_addr, client_addr;
    socklen_t addr_len = sizeof(client_addr);

    server_socket = socket(AF_INET, SOCK_STREAM, 0);
    if (server_socket == -1) {
        cerr << "Failed to create socket.\n";
        return 1;
    }

    server_addr.sin_family = AF_INET;
    server_addr.sin_addr.s_addr = INADDR_ANY;
    server_addr.sin_port = htons(PORT);

    if (bind(server_socket, (struct sockaddr*)&server_addr, sizeof(server_addr)) < 0) {
        cerr << "Bind failed.\n";
        return 1;
    }

    if (listen(server_socket, 5) < 0) {
        cerr << "Listen failed.\n";
        return 1;
    }

    cout << "Server is listening on port " << PORT << "...\n";

    // Thread for accepting new clients
    thread([&]() {
        while (true) {
            client_socket = accept(server_socket, (struct sockaddr*)&client_addr, &addr_len);
            if (client_socket < 0) {
                cerr << "Accept failed.\n";
                continue;
            }
            cout << "New client connected.\n";
            client_sockets.push_back(client_socket);

            // Handle each client in a separate thread
            thread(handle_client, client_socket).detach();
        }
    }).detach();

    // Main thread for sending messages from the server
    string server_message;
    while (true) {
        getline(cin, server_message);
        if (server_message == "exit") {
            break;
        }
        broadcast_message("Server: " + server_message);
    }

    close(server_socket);
    return 0;
}


\end{lstlisting}

\subsection*{Client Code}
\begin{lstlisting}[language=C++]
#include <iostream>
#include <thread>
#include <cstring>
#include <sys/socket.h>
#include <arpa/inet.h>
#include <unistd.h>
using namespace std;

#define PORT 8094
#define BUFFER_SIZE 1024

void receive_messages(int socket) {
    char buffer[BUFFER_SIZE];
    while (true) {
        memset(buffer, 0, BUFFER_SIZE);
        int bytes_read = recv(socket, buffer, BUFFER_SIZE, 0);
        if (bytes_read <= 0) {
        cerr << "Disconnected from server or error receiving message.\n";
            break;
        }
        cout << "Broadcast message received: " << buffer << "\n";
    }
}

int main() {
    int client_socket;
    struct sockaddr_in server_addr;

    client_socket = socket(AF_INET, SOCK_STREAM, 0);
    if (client_socket == -1) {
        cerr << "Failed to create socket.\n";
        return 1;
    }

    server_addr.sin_family = AF_INET;
    server_addr.sin_port = htons(PORT);
    server_addr.sin_addr.s_addr = inet_addr("127.0.0.1");

    if (connect(client_socket, (struct sockaddr*)&server_addr, sizeof(server_addr)) < 0) {
        cerr << "Connection failed.\n";
        return 1;
    }

    cout << "Connected to server.\n";

    // Start a thread to receive messages
    thread(receive_messages, client_socket).detach();

    // Main loop for sending messages
    while (true) {
        string message;
        getline(cin, message);
        if (message == "exit") {
            break;
        }
        send(client_socket, message.c_str(), message.size(), 0);
    }

    close(client_socket);
    return 0;
}

\end{lstlisting}
\newpage

\begin{figure}
    \centering
    \includegraphics[width=1\linewidth]{6.png}
    \label{fig:enter-label}
\end{figure}


\section*{Part 2: Distributed System Socket Programming}

\subsection{Server Code}
\begin{lstlisting}
#include <iostream>
#include <unistd.h>
#include <string.h>
#include <sys/socket.h>
#include <arpa/inet.h>
using namespace std;

#define PORT 8084

int main()
{
    int server_fd, client_sock;
    struct sockaddr_in server_addr, client_addr;
    socklen_t client_len = sizeof(client_addr);
    char buffer[1024] = {0};

    // Create socket
    if ((server_fd = socket(AF_INET, SOCK_STREAM, 0)) == 0)
    {
        perror("Socket failed");
        exit(EXIT_FAILURE);
    }

    server_addr.sin_family = AF_INET;
    server_addr.sin_addr.s_addr = INADDR_ANY;
    server_addr.sin_port = htons(PORT);

    // Bind the socket to the port
    if (bind(server_fd, (struct sockaddr *)&server_addr, sizeof(server_addr)) < 0)
    {
        perror("Bind failed");
        exit(EXIT_FAILURE);
    }

    // Listen for incoming connections
    if (listen(server_fd, 3) < 0)
    {
        perror("Listen failed");
        exit(EXIT_FAILURE);
    }
    cout << "Waiting for connections...\n";

    // Accept a connection from the client
    if ((client_sock = accept(server_fd, (struct sockaddr *)&client_addr, &client_len)) < 0)
    {
        perror("Accept failed");
        exit(EXIT_FAILURE);
    }

    cout << "Client connected\n";

    // Read and send messages
    while (true)
    {
        int read_size = read(client_sock, buffer, 1024);
        if (read_size == 0)
        {
            cout << "Client disconnected\n";
            break;
        }
        buffer[read_size] = '\0';
        cout << "Client: " << buffer << "\n";

        cout << "Server: ";
        string response;
        getline(cin, response);
        send(client_sock, response.c_str(), response.length(), 0);
    }

    close(client_sock);
    close(server_fd);           
    return 0;
}



\end{lstlisting}
\begin{figure}
    \centering
    \includegraphics[width=1\linewidth]{distributed.png}
    \caption{distributed server}
    
\end{figure}
\newpage
\subsection{Client Code}
\begin{lstlisting}
#include <iostream>
#include <unistd.h>
#include <string.h>
#include <sys/socket.h>
#include <arpa/inet.h>
using namespace std;

#define SERVER_IP "127.0.0.1" // Replace with server laptop's IP address
#define PORT 8080

int main()
{
    int sock;
    struct sockaddr_in server_addr;
    char buffer[1024] = {0};

    // Create socket
    if ((sock = socket(AF_INET, SOCK_STREAM, 0)) < 0)
    {
        perror("Socket creation failed");
        exit(EXIT_FAILURE);
    }

    server_addr.sin_family = AF_INET;
    server_addr.sin_port = htons(PORT);

    // Convert IP address to binary format
    if (inet_pton(AF_INET, SERVER_IP, &server_addr.sin_addr) <= 0)
    {
        perror("Invalid address or Address not supported");
        exit(EXIT_FAILURE);
    }

    // Connect to the server
    if (connect(sock, (struct sockaddr *)&server_addr, sizeof(server_addr)) < 0)
    {
        perror("Connection failed");
        exit(EXIT_FAILURE);
    }

    cout << "Connected to server\n";

    // Send and receive messages
    while (true)
    {
        cout << "Client: ";
        string message;
        getline(cin, message);
        send(sock, message.c_str(), message.length(), 0);

        int read_size = read(sock, buffer, 1024);
        buffer[read_size] = '\0';
        cout << "Server: " << buffer << "\n";
    }

    close(sock);
    return 0;
}


\end{lstlisting}

\newpage
\begin{figure}
    \centering
    \includegraphics[width=1\linewidth]{another.png}
    \caption{distributed client}
   
\end{figure}
\end{document}
